\section{{Methodology and Experimental Design}}
\label{sec:method}

\subsection{Applicability conditions and how a pipeline can lose them}

A question about a policy often has an answer that holds only under a condition the asker has not settled. A person asks whether they can bring another child on a Parent of a Child Student visa; the guidance says yes, \emph{if} the child also has or is applying for a Child Student visa. Neither ``yes'' nor ``no'' is a correct answer on its own: the correct answer is ``yes, provided that \dots''. We call the clause that must also hold an \emph{applicability condition}.

Formally, a question $q$ with scenario $s$ has a gold answer set $A(q)=\{(a_i, C_i)\}$, where $a_i$ is an answer and $C_i$ is the set of condition statements under which $a_i$ holds and which $s$ leaves unresolved. A question is \emph{conditional} when some $C_i \neq \emptyset$. The statements in $C_i$, together with the statements that give $a_i$ itself (its \emph{answer evidence} $E_i$), are passages of a source document $d(q)$.

A retrieval-augmented pipeline never sees $d(q)$ whole. It sees chunks. Let $\kappa$ be a chunking function that partitions a document into chunks, and let $\kappa(x)$ denote the chunk holding statement $x$. For each conditional question we define
\begin{equation}
\mathrm{Sep}(q) = \mathbf{1}\Big[\big|\{\kappa(x) : x \in \textstyle\bigcup_i (C_i \cup E_i)\}\big| > 1\Big],
\end{equation}
the event that the gold no longer fits in one chunk, and the set of \emph{condition-only chunks}
\begin{equation}
\mathcal{K}_C(q) = \{\kappa(x) : x \in \textstyle\bigcup_i C_i\} \setminus \{\kappa(x) : x \in \textstyle\bigcup_i E_i\},
\end{equation}
chunks that carry a condition and nothing else the answer needs. A chunk in $\mathcal{K}_C(q)$ is exactly the kind of passage a pipeline can drop while still producing a fluent, confident and partly correct answer.

Applicability conditions take several logical forms: eligibility requirements (``if you are over 18''), exceptions (``unless the ATF repairs the vehicle''), time limits, prohibitions and conditional permissions. ConditionalQA's annotators mark as a condition any statement an answer depends on that the scenario leaves open, whatever its form, and we follow their annotation rather than impose a finer taxonomy; the forms are not distinguished in any result below. Conditional evaluation questions carry \annMeanCond{} condition elements on average, \annTwoPlus\% have two or more, and \annMultiAnswer\% have more than one answer, each with its own conditions.

A condition can be lost at four stages, which this work separates:
\begin{enumerate}[(i)]
\item \textbf{Chunking} places it in a chunk apart from the answer.
\item \textbf{Retrieval} ranks that chunk below the cut-off $k$.
\item \textbf{Evidence reduction}: a reranker, relevance grader, contradiction filter or conflict resolver removes the chunk after retrieval.
\item \textbf{Generation}: the condition is in the context and the model does not state it, or states a condition that does not apply.
\end{enumerate}

\subsection{Testbed: ConditionalQA as chunked open retrieval}

ConditionalQA v1.0 \cite{sun2022conditionalqa} is built on \nPages{} UK government guidance pages from gov.uk; every answer is annotated with the conditions under which it holds, and every condition and evidence statement is quoted verbatim as an HTML element of its page (a paragraph, list item, table row or heading). We turn it into open retrieval over the whole collection.

\textbf{Chunking.} Each page is cut into chunks of whole elements, packed greedily up to 150 words, so the chunk holding every gold statement is known exactly and $\mathrm{Sep}(q)$ and $\mathcal{K}_C(q)$ follow directly from the annotation: \nChunks{} chunks, mean \chunkWords{} words, each prefixed for retrieval with its page title and nearest section heading.

\textbf{Questions.} The test answers are not public, so we evaluate on the labelled training and development questions; all systems are zero-shot, so both are unseen to them. Before running any system we fixed a document-disjoint split: every question on \nTuningPages{} random pages forms a \emph{tuning} set (\nTuning{} questions), the only place prompts and thresholds were chosen; the \emph{evaluation} set holds every conditional (\nEvalCond{}) and unanswerable (\nEvalUnans{}) question on the other pages and \nEvalUncond{} random unconditional ones, \nEval{} questions on \nEvalPages{} pages.

\textbf{Retrieval.} BM25 \cite{robertson2009bm25} and Contriever \cite{izacard2022contriever} rankings are fused by reciprocal rank \cite{cormack2009rrf} over all chunks; the query is the question followed by the scenario, and the rankings are shared by every system.

\subsection{One generator, many evidence policies}

Every system is defined by one thing, \emph{which chunks reach the generator}; the generator, its instructions and its output format are identical for all, so any difference between two systems is a difference in their evidence. The generator sees the chunks as numbered lines, one per HTML element, with the scenario and question, and returns JSON listing each answer with the lines that state it (\emph{evidence}) and the lines that state a requirement which must also hold and which the scenario does not settle (\emph{conditions}); condition lines are mapped back to their exact elements, the form the official metric scores. Four worked examples from the tuning pages precede the question. Separating evidence from conditions was the one change the tuning set showed to matter (Section~\ref{sec:generation}). Because citations are line numbers, a condition cannot be missed through paraphrase: a gold condition in context is cited as a condition, cited as evidence, or not cited.

\textbf{Generation settings.} Decoding is greedy (temperature 0). Lines are numbered in page order, identically for every model, and each model receives the same messages through its provider's chat template. The output budget is 768 tokens for the local models, 3{,}072 for gpt-oss-20b (low reasoning effort) and 1{,}024 for Gemma-4-31B and Nemotron-3-Ultra; the local context window is 8{,}192 tokens; at most five answers are scored. An output that is not valid JSON in the required shape, is cut off at the token limit, or is aborted by the server as a repetition loop counts as unanswered; provider errors are retried, and the recorded answer is the model's first complete output. Prompts, examples and settings are in the repository.

Table~1 lists the systems, in three families. \textbf{The oracle ladder} isolates stages by construction: \texttt{gold\,$-$\,cond.\ elements} measures how much the metric depends on conditions being present at all; \texttt{gold\,$-$\,cond.-only chunks} is what a pipeline delivers when it finds the answer and loses a separately chunked condition; \texttt{gold\,+\,distractors} gives perfect recall under realistic distraction; a closed-book run measures what the models answer from memory. \textbf{Evidence-reducing operators} start from the top 8 and reproduce modules common in robust RAG stacks: a cross-encoder reranker (checkpoint \texttt{BAAI/bge-reranker-base}, BGE family \cite{xiao2024cpack}); a contradiction filter built on an NLI cross-encoder \cite{he2023debertav3} (an NLI judgement says only that two passages cannot both hold as stated, not which is authoritative or whether one qualifies the other, which is precisely the case at issue); a ConflictRAG-style resolver \cite{conflictrag2026}; relevance grading in the manner of corrective RAG \cite{yan2024crag}; and listwise selection, a single-call reduction of Astute RAG's consolidation \cite{wang2025astute}. The resolver is a partial reproduction, since the authors' code is not released: it scores three of the original five credibility criteria (source authority, recency and specificity), omitting the relevance and consistency criteria the original extracts with a language model, and on gov.uk pages authority and recency tie, so specificity decides; our results concern this implementation, not the published system. \textbf{Context completion}, described next, adds chunks instead; its budget-matched baselines add chunks without any knowledge of conditions.

\begin{table}[!htbp]\centering
\textbf{Table 2.} The systems. Each differs only in which chunks reach the same generator. Operators are marked by how closely they follow the method they are named after: [S] a standard component, not a reproduction; [P] a partial reproduction; [R] a reduction of the method to its evidence-reducing step. $^\dagger$\,Exploratory in the analysis plan; $^\ddagger$\,added in revision (exploratory, except the clean threshold protocol, a corrected re-analysis that is primary for completion).\\\vspace{5pt}
{\fontsize{9}{11}\selectfont
\begin{tabular}{L{4.3cm}L{10.6cm}}
\toprule
\textbf{System} & \textbf{Chunks given to the generator} \\
\midrule
\multicolumn{2}{l}{\emph{Oracle ladder (uses the gold annotation; not deployable)}}\\
\texttt{closed-book} & none \\
\texttt{gold} & exactly the chunks holding the annotated evidence and conditions \\
\texttt{gold\,$-$\,cond.\ elements} & the same chunks with the condition elements deleted \\
\texttt{gold\,$-$\,cond.-only chunks} & the gold chunks minus those holding only conditions \\
\texttt{gold\,+\,distractors} & all gold chunks, filled to 8 with top-ranked non-gold chunks \\
\midrule
\multicolumn{2}{l}{\emph{Retrieval}}\\
\texttt{top-$k$} & the $k$ highest-ranked chunks, $k\in\{3,8\}$ (recall alone is measured up to $k=50$); \texttt{top-8} is the reference \\
\midrule
\multicolumn{2}{l}{\emph{Evidence-reducing operators (applied to the top 8)}}\\
\texttt{rerank-3} & [S] the 3 chunks a cross-encoder reranker scores highest \\
\texttt{NLI filter} & [S] drops the lower-ranked chunk of each pair judged contradictory \\
\texttt{ConflictRAG-style} & [P] resolves each detected conflict by Entropy-TOPSIS credibility selection (3 of 5 criteria) \\
\texttt{relevance grading}$^\dagger$ & [R] corrective RAG's grading: keeps the chunks the generator model grades as useful, one call each (RAMDocs only) \\
\texttt{listwise selection}$^\dagger$ & [R] Astute RAG's consolidation in one call: keeps the chunks the generator model selects \\
\midrule
\multicolumn{2}{l}{\emph{Context completion (budget of 8 chunks)}}\\
\texttt{adjacent}$^\dagger$ & the top 5 plus the chunks next to them on their pages \\
\texttt{page}$^\dagger$ & every chunk of the top-ranked chunk's page \\
\texttt{completion (untrained)}$^\dagger$ & the top 5 plus same-page chunks scored by an off-the-shelf reranker \\
\texttt{completion} (ours)$^\ddagger$ & the top 5 plus same-page chunks a trained scorer predicts state a condition; each fold's scorer and threshold set without any evaluation label \\
\texttt{completion (pre-reg.\ $\tau$)} & the same, with the pre-registered threshold, chosen under a scorer trained on the evaluation folds \\
\midrule
\multicolumn{2}{l}{\emph{Context-preserving retrieval$^\ddagger$ (budget: the words of the top 8)}}\\
\texttt{sentence window} & single elements retrieved, each passed with two neighbours on either side \\
\texttt{parent-child} & single elements retrieved, each passed with its top-level section \\
\texttt{hierarchical} & the top 2 pages retrieved first, then their 8 highest-ranked chunks \\
\texttt{query expansion} & the top 8 chunks for the query extended with condition cue words \\
\bottomrule
\end{tabular}
}
\end{table}


\subsection{Qualifier-aware context completion}

We test one remedy aimed at the retrieval side, stated as a hypothesis: if the retriever has reached the right page but not the chunk that qualifies the answer, a scorer trained on the benchmark's condition annotation can find that chunk among its neighbours. By design it can only add chunks from pages the top chunks already reached; Section~\ref{sec:completion} measures how often the missing condition lies there. Qualifier-aware completion keeps the top $k_0=5$ retrieved chunks and adds, from the pages those chunks come from, the chunks most likely to state a condition on the answer, up to the same budget of eight chunks as \texttt{top-8}.

For a question $q$ with scenario $s$, let $P$ be the (at most two) pages of the top-$k_0$ chunks and, for each page $p\in P$, let $a_p$ be its highest-ranked retrieved chunk, the \emph{anchor}. Every other chunk $c$ on $p$ that was not retrieved is a candidate, scored by a cross-encoder $f_\theta$:
\begin{equation}
\sigma(c) = f_\theta\big([\,q;\, s;\, a_p\,],\; c\big) \in [0,1].
\end{equation}
The system adds the candidates with $\sigma(c)\geq\tau$, highest first, up to $8-k_0$ of them, and leaves the remaining slots empty: it adds qualifiers, not context for its own sake.

$f_\theta$ is DeBERTa-v3-base \cite{he2023debertav3} fine-tuned to predict whether a candidate chunk holds a gold condition element of the question. The labels are ConditionalQA's own annotation, so no new labelling is needed; unconditional questions contribute negatives only, which is what teaches the scorer when to add nothing. To keep every question scored by a model that never saw its page, training is \emph{cross-fitted}: the evaluation questions are split into five document-disjoint folds, and the model that scores a fold is trained on the other four folds and the tuning set. The tuning questions are scored by a model trained on all five evaluation folds, and $\tau$ is chosen on them alone. That model has seen evaluation labels, but it never scores an evaluation question; it only sets $\tau$, and Section~\ref{sec:completion} shows that no threshold, even one chosen on the evaluation questions themselves, makes the scorer accurate, so the choice of $\tau$ cannot account for the result. A second variant, \texttt{completion (untrained)}, scores candidates with the off-the-shelf reranker on $(q, c)$ instead, to separate what training contributes from what any relevance model would give. A third, \texttt{completion (oracle)}, added in revision, adds exactly the candidates that hold a gold condition: it shows what the architecture could gain with a perfect scorer.

\subsection{A secondary diagnostic: the contrast regime}

ConditionalQA is the primary evaluation. RAMDocs \cite{wang2025madam} serves only as a secondary diagnostic, and its limits come first: its misinformation is synthetic, \nRamMulti{} of its \nRamQ{} questions have several valid answers for same-name entities, so it mixes ambiguity with conflict, and it has no training split. Its absolute scores are therefore not comparable with ConditionalQA's, and we draw from it only one contrast. Removing evidence is not always wrong: when some retrieved passages are false, discarding them is the purpose of the operators above. Because each RAMDocs document is labelled correct, misinformation or noise, it allows an \emph{oracle discard} system that removes exactly the misinformation and noise, the ceiling for any operator that removes evidence, and so asks whether the same operators that cannot tell a condition from a contradiction can at least tell a false document from a true one. A RAMDocs answer is correct when it contains every gold answer and no wrong answer.

\subsection{Metrics and inference}

The primary metric is ConditionalQA's official \emph{F1 with conditions}, computed by the authors' evaluation script, unchanged. For one gold answer $(a, C)$ and one predicted answer $(\hat a, \hat C)$ it is
\begin{equation}
\mathrm{F1}_c = \mathrm{F1}_{\mathrm{tok}}(\hat a, a)\cdot
\begin{cases}
\mathbf{1}[\hat C=\emptyset] & \text{if } C=\emptyset,\\
\mathrm{F1}_{\mathrm{set}}(\hat C, C) & \text{otherwise,}
\end{cases}
\end{equation}
where $\mathrm{F1}_{\mathrm{tok}}$ is SQuAD token F1 and $\mathrm{F1}_{\mathrm{set}}$ is the F1 of exact element matches. Over several answers the script takes the best one-to-one assignment and multiplies by $\exp(1-|\hat A|/|A|)$ when more answers are predicted than exist. The metric therefore cannot be raised by verbosity: listing extra answers is discounted, and naming a condition where the gold has none scores zero. Because this score rewards a generator that never states a condition on unconditional questions, we never read it alone. We also report answer F1 (conditions ignored), the precision and recall of predicted condition elements, exact match with conditions (answer and conditions both exactly right), the share of unconditional answers that name any condition (false conditions), how the gold conditions in context are handled by the generator (stated as conditions, cited only as evidence, or not cited), and, for unanswerable questions, the share left unanswered.

Questions on the same page share chunks and retrieval behaviour, so they are not independent. Every interval and test treats the \emph{page} as the unit: differences between two systems are paired by question, with 95\% BCa bootstrap intervals over pages (2{,}000 resamples) and page-level sign-flip permutation tests (20{,}000 permutations). The confirmatory comparisons, listed in the analysis plan we committed before the first evaluation run, are corrected by the Holm procedure \cite{holm1979} separately for each model.

\textbf{Models.} The two primary generators run locally on a consumer GPU: Qwen2.5-7B-Instruct \cite{qwen2024qwen25} and Llama-3.1-8B-Instruct \cite{grattafiori2024llama3}, both 4-bit quantised, at temperature 0. To test whether the direction of the main effects holds at larger scale, we repeat the key systems on a stratified subsample of \nHostedQ{} evaluation questions with three larger open-weight models from three families: gpt-oss-20b \cite{openai2025gptoss}, a mixture of experts served by Groq and, once Groq's daily quota is spent, by NVIDIA's API with identical settings (Section~\ref{sec:hosted} checks the two providers against each other on the same questions); Gemma-4-31B (NVIDIA API identifier \texttt{google/gemma-4-31b-it}; the checkpoint published as \texttt{google/gemma-4-31B-it}), a dense model; and Nemotron-3-Ultra (\texttt{nvidia/nemotron-3-ultra-550b-a55b}), a hybrid mixture of experts with 55B of its 550B parameters active, run with its reasoning mode off, both served by NVIDIA. Every answer records the provider that served it. These are reported as replication, not as confirmatory tests. The models differ in architecture, training, quantisation and serving as well as in size, so their comparison is not a controlled scaling experiment. Qwen3.8-27B (Groq model identifier \texttt{qwen/qwen3.8-27b}) was used only in prompt tuning (Section~\ref{sec:generation}).

\textbf{Diagnostics added in revision.} Four exploratory analyses were added after the confirmatory results, and are recorded as such in the analysis plan. (i) A \emph{condition-verification pass}: after the frozen generator has answered, a second call to the same model receives the same lines, the scenario, the question and the answers, and lists for each answer the lines that state a requirement the scenario does not settle; the answers are kept and only their conditions replaced. It tests whether a dedicated check recovers the conditions the generator ignores. (ii) \texttt{completion (oracle)}, described above. (iii) The \texttt{top-8} and \texttt{gold} systems with section-aware chunks instead of 150-word chunks, end to end, with a fixed 70-word chunking whose mean size matches the section-aware chunks as a control. (iv) Every gpt-oss-20b gold-context question that Groq served is asked again on NVIDIA's API.

\textbf{Transfer corpus.} The out-of-distribution check of the completion scorer uses a corpus from the authors' earlier work: \nOodCorpus{} query-and-passage-pair instances from gov.uk immigration guidance and UK bank terms, labelled for conflict type and scope relation (a refinement is a pair in which one passage states a rule and the other a condition or exception to it), including written probe pairs and distractor pairs with no conflict. Each instance carries two annotators' labels, but no valid inter-annotator agreement was established for it, so its labels are treated as provisional and the check as indicative only.
